\documentclass[ekobook.tex]{subfiles}

\begin{document}\setWPP
\setlength{\multicolsep}{0.5em}% Remove vertical space before and after
\setlength{\columnsep}{0.5em}% Remove space between columns
\setlength{\leftskip}{0em}%
\lsection{Management}
\qaw{Pojem management}{%
\texthl{Soubor vedoucích pracovníků}, kteří svojí činností \texthl{zajišťují prosperitu podniku}.
}
\qaw{Skupiny manažerů, dovednosti}{%
\begin{multicols}{2}
\begin{wnumbered}
\wtem{ TOP management}
zodpovídá za \texthl{nejdůležitější strategická rozhodnutí} (ředitel/ka)
\wtem{střední management}
zodpovídá za \texthl{své úseky firmy} (zástupce ředitele, vedoucí ekonomického úseku, ...)
\wtem{nižší management}
každodenní \texthl{řešení operativních situací} (mistr na směně, ...)
\end{wnumbered}
\columnbreak
\begin{enumerate}[label={\color{WPP}\textbf{\Alph*)}}]
\wtem{Naučené dovednosti (hard skills)}
dovednosti získané \texthl{během studia a praxe} (technická odbornost, cizí jazyk, ...)
\wtem{Vrozené dovednosti (soft skills)}
dovednosti \texthl{jednat s lidmi} (komunikace, schopnost vyjednávat, empatie, ...)
\end{enumerate}
\end{multicols}
}
\qaw{Manažerské funkce}{%
\begin{enumerate}[label={\color{WPP}\texthl{\alph*.}}]
\item \texthl{plánování}
\item \texthl{organizování}
\item \texthl{vedení}
\item \texthl{motivace}
\item \texthl{kontrola}
\end{enumerate}
}
\qaw{Druhy plánů}{%
\textbl{Dlouhodobé (strategické)}
-- nad pět let; například stavba haly\\
\textbl{Střednědobé (taktické)}
-- nad rok, do pěti let; nákup strojů, rozvržení výroby, ...\\
\textbl{Krátkodobé (operativní)}
-- do jednoho roku; rozhodování každodenných činností, plán výroby na příští měsíc, ...
}
\qaw{Vysvětli funkcionální a lineární organizaci podniku}{%
\textbl{Funkcionální organizace} udává několik manažerů, každého schopného ve svém odvětví.
Zaměstnanec \texthl{musí vyhledat vedoucího} dle toho, co potřebuje.
Chce vědět, co má ten den dělat? Jde za šéfem výroby.
Zajímá jej škrt na výplatní pásce? Jde za šéfem účetní jednotky.
\texthl{Vhodná pro menší firmy.}\\
\textbl{Lineární organizace} udává každému pracovníkovi \texthl{jednoho přímého nadřízeného}.
Každý nadřízený má svůj \texthl{pevně daný tým}, kterému věnuje svůj čas.
\texthl{Vhodné pro větší firmy kvůli přehlednosti.}\\
Největší firmy často používají \textbl{smíšený model}.
Lineární organizace pro každodenní pokyny, ale oddělení specialistů, co také mají pravomoce nad běžnými pracovníky.
}
\qaw{Druhy motivace}{%
\begin{multicols}{2}
\textbl{pozitivní hmotná}
-- prémie; služební auto; stravenky; ...\\
\textbl{pozitivní nehmotná}
-- povýšení\\
\textbl{negativní hmotná}
-- snížení mzdy; odebrání auta; ...\\
\textbl{negativní nehmotná}
-- napomenutí; sesazení z funkce; ...
\end{multicols}
}
\qaw{Popiš teorie vedení a styly řízení}{%
\begin{multicols}{2}
\textbl{Autokratický styl}
-- vedoucí vydává \texthl{absolutní příkazy} k okamžitému splnění, hojně využívá \texthl{teorie X}; výhoda spočívá v možnosti \texthl{rychlého rozhodování} a řešení vzniklých situací\\
\textbl{Demokratický styl}
-- vedoucí si ponechává své \texthl{právo posledního slova}, ale častěji diskutuje chod firmy se zaměstnanci, nabízí jim možnost \texthl{teorie Y} a oceňuje vstup zaměstnanců\\
\textbl{Liberální styl}
-- vedoucí využívá pouze \texthl{nepřímých příkazů} a nechává tým pracovat sám o sobě; vhodný například u výzkumných týmů, operačních sálů; kde každý ví, co je jeho prací a je motivován ji plnit
\columnbreak\\
\textbl{Teorie X (krátké vodítko)}
-- očekává se \texthl{pouze přesné plnění příkazů}; zaměstnanec pracuje jako ozubené kolečko ve stroji bez vkládání vlastní mysli; použitelný u neměnné manuální práce, jako je práce u pásu\\
\textbl{Teorie Y (dlouhé vodítko)}
-- od zaměstnance \texthl{se očekává aktivita}, seberealizace, zapojení se do chodu firmy, vklad vlastních nápadů
\end{multicols}
}
\qaw{Jaká může být vnější a vnitřní kontrola}{%
\textbl{Vnější kontrola} může být \texthl{zákonná} nebo \texthl{smluvní}. Patří sem hygiena, ČOI, finanční úřad. Smluvní se pak využívá třeba při žádosti o úvěr, kdy banka dostane možnost nahlédnout do chodu firmy před schválením žádosti.\\
\textbl{Vnitřní} je pak v rámci firmy a vykonávají ji \texthl{specializovaní pracovníci}.\\
Kontroly se dělí na plánované \wx{ } neplánované; částečné \wx{ } celkové; mimořádné \wx{ } pravidelné.
}
\wpage
\end{document}